\documentclass{beamer}
\usetheme{STCE}

\usepackage{beamerfoils}
\usepackage[utf8]{inputenc}
\usepackage{amsmath}
\usepackage[german]{babel} % Für deutche Umlaute
\usepackage{tabularx} % Für die TabularX Umgebung

\MyLogo{
\centering
\includegraphics[width=0.12\textwidth]{./figures/logo.eps}
}

\begin{document}
\title{Correlation and Regression Analysis}
\author{Patrick Neidig \and Marius Grysla}
\institute{Software and Tools for Computational Engeneering}
\date{31.5.2011}
\frame{\titlepage}

\begin{frame}
 \frametitle{Inhalt}
 \tableofcontents
\end{frame}

\section{Korrelation}
\begin{frame}
 
\end{frame}


\section{Multiple Regression}
\begin{frame}
 \frametitle{Das multiple lineare Regressionsmodell}
 \begin{itemize}
 \item 
 \end{itemize}
\end{frame}


\begin{frame}
 \frametitle{Das multiple lineare Regressionsmodell}

 \begin{itemize}
  \item Wie groß sind die Einflüsse von verschiedenen Regressanden auf ein abhängiges Merkmal?
 \end{itemize}
 
 \begin{block}{Lineare Einfachregression}
  \begin{equation*}
   Y_i = \alpha + \beta x_i + \epsilon_i, \quad i = 1, \dots, n
  \end{equation*}
 \end{block}

 \begin{block}{Multiple Lineare Regression}
  \begin{equation*}
   Y_i = \beta_0 + \beta_1 x_{i1} + \dots + \beta_p x_{ip} + \epsilon_i, \quad i = 1, \dots, n
  \end{equation*}
  \begin{itemize}
   \item Anzahl der Regressanden: $p$
  \end{itemize}

 \end{block}

\end{frame}

\begin{frame}
  \frametitle{Anwendung}
  % Hier sollte die Verwendung der Regression mit dem Mietspiegel Beispiel verdeutlicht werden
  \begin{itemize}
  \item Beispieldatensatz:\\
    \begin{tabular}[h]{c|c|c|c}
      1 & 2 & 3 & 4 \\
      5 & 6 & 7 & 8
    \end{tabular}
  \end{itemize}

\end{frame}

\begin{frame}
  \frametitle{Funktionssignatur}
  
  \begin{itemize}
  \item Funktionsname: nag\_regsn\_mult\_linear (g02dac)
  \item Argumente:
    \begin{itemize}
    \item Anzahl der Beobachtungen $n$
    \item $m \times n$ Matrix $X$ mit den Beobachteten Werten für alle Variablen
    \item \dots
    \end{itemize}
  \end{itemize}

\end{frame}

\subsection{NAG Algorithmus}
\begin{frame}
  \frametitle{NAG Algorithmus}
 
  \begin{block}{Für Berechnung wird eine Matrixschreibweise benötigt}
    {\centering $y = X \beta + \epsilon$ \\}
    wobei \\
    \qquad $y$ = Vektor mit $n$ Beobachtungen von $y$ \\
    \qquad $X$ = Matrix mit den Beobachtungen der Regressanden \\
    \qquad $\beta$ = Vektor mit unbekannten Faktoren \\
    \qquad $\epsilon$ = Vektor mit Fehlertermen \\
  \end{block}
  
  \pause
  
  \begin{block}{Methode der kleinsten Quadrate}
    In Matrixschreibweise:\\
    {\centering
      $\sum\limits^{N}_{n=1} \epsilon^2_n = \epsilon^T \epsilon = (y - X \beta)^T (y - X \beta) \rightarrow \min\limits_{\beta}$\\}
    oder: \\
    {\centering
      $\Vert X\beta - y \Vert_2 \rightarrow \min\limits_{\beta}$
      \\}
  \end{block}

\end{frame}



\begin{frame}
  \frametitle{QR-Zerlegung}
  
  \begin{itemize}
  \item Wie kann das benötigte Minimum berechnet werden?\\
    $\Rightarrow$ QR-Zerlegung der Matrix $X$
  \end{itemize}

  \begin{block}{QR-Zerlegung}
    Zerlegung von $X$ in zwei Matrizen $Q$ und $R$, sodass \\
    \qquad $Q$ ist orthogonale Matrix \\
    \qquad $R$ enthält obere Dreiecksmatrix \\
  \end{block}

  \pause

  \begin{itemize}
  \item Algorithmen:
    % Hier Vor-/Nachteile erwähnen
    \begin{itemize}
    \item Householder Transformationen 
    \item Givens Rotationen 
    \item Modifizierter Gram-Schmidt Algorithmus 
    \end{itemize}
  \end{itemize}

\end{frame}

\begin{frame}
  \frametitle{QR-Zerlegung (Beispiel)}
  % Eine QR-Zerlegung mit den Beispieldaten durchführen
  % Zum Beispiel die ersten Schritte der Givens Rotation durchführen
\end{frame}

\begin{frame}
  \frametitle{Gleichungssystem Lösen}
  
  \begin{itemize}
  \item Wir haben: $X = QR^*$, wobei $R^* = \binom{R}{0}$
  \item Es gilt: $\Vert X\beta - y \Vert_2 = \Vert Q^T X \beta - Q^T y \Vert_2 = \Vert R \beta - c \Vert_2 + \Vert d \Vert_2w$
  \item $Q^T * y = \binom{c}{d}$
  \item Löse $R \hat{\beta} = c$ durch Rückwärtseinsetzen
    
  \pause

  \item Funktioniert nur bei vollem Rang von $R$!
  \end{itemize}

\end{frame}

\begin{frame}
  \frametitle{Gleichungssystem Lösen (Beispiel)}
  
\end{frame}

% Zu Komplex für den Vortrag?
%\begin{frame}
%  \frametitle{Gleichungssystem Lösen II}
%  
%\end{frame}

\begin{frame}
  \frametitle{Ergebnis}
  % Plots mit Mieten im Verhältnis zu Regressoren zeigen
  % Ergebnis der Beispielrechnung mit dem Regressionsergebnis mit vollständigen Daten vergleichen. 
\end{frame}

\begin{frame}
  \frametitle{Ausblick}
  
  Im nächsten Vortrag folgt:
  \begin{itemize}
  \item Wie werden die Methoden angewendet
    \begin{itemize}
    \item Am Beispiel des Münchener Mietspiegels
    \end{itemize}
  \item Performance Analyse
    \begin{itemize}
    \item Wie verhalten sich die Algorithmen bei verschiedenen Datensätzen/Parametern 
    \end{itemize}
  \end{itemize}

\end{frame}

\end{document}
